\documentclass[preprint,aps,nofootinbib,a4paper,titlepage,superscriptaddress,longbibliography,amsfonts,amssymb,amsmath]{revtex4-2}

\newcommand{\mr}{\ms\hline\ms}
\usepackage{etoolbox}
\usepackage{IEEEtrantools}
\usepackage{microtype}
\usepackage[toc,title,page,titletoc]{appendix}
\usepackage{amsmath}
\usepackage{amsthm}
\usepackage{mathrsfs}
\usepackage{enumitem}
\usepackage[T1]{fontenc}
\usepackage{xfrac}
\usepackage{lmodern}
\usepackage[a4paper]{geometry}
\usepackage{hyperref}
\usepackage{mathbbol}

%Tikz
\usepackage{tikz}
\usetikzlibrary{patterns}
\usetikzlibrary{shapes.geometric}
\usetikzlibrary{decorations.pathmorphing}
\usetikzlibrary{arrows.meta}
\tikzset{snake it/.style={decorate, decoration=snake}}
\usetikzlibrary{positioning}
\usetikzlibrary{arrows,shapes,positioning}
\usetikzlibrary{decorations.markings}
\tikzstyle arrowstyle=[scale=1]
\tikzstyle directed=[postaction={decorate,decoration={markings,mark=at position .65 with {\arrow[arrowstyle]{stealth}}}}]
\tikzstyle reverse directed=[postaction={decorate,decoration={markings,mark=at position .65 with {\arrowreversed[arrowstyle]{stealth};}}}]

\tikzset{->-/.style={decoration={
  markings,
  mark=at position #1 with {\arrow{>}}},postaction={decorate}}}

\tikzset{-<-/.style={decoration={
  markings,
  mark=at position #1 with {\arrow{<}}},postaction={decorate}}}

% Calligraphic letters
\renewcommand{\S}{\mathcal{S}}
\newcommand{\E}{\mathcal{E}}
\newcommand{\T}{\mathcal{T}}
% \newcommand{\Lie}{\mathcal{L}}
\renewcommand{\O}{\mathcal{O}}
\renewcommand{\P}{\mathcal{P}}

\newcommand{\sM}{\mathscr{M}}
\newcommand{\sS}{\mathscr{S}}
\newcommand{\sI}{\mathscr{I}}
\newcommand{\sH}{\mathscr{H}}
\newcommand{\sP}{\mathscr{P}}
\newcommand{\sL}{\mathscr{L}}
\newcommand{\sV}{\mathscr{V}}
\newcommand{\sW}{\mathscr{W}}
\newcommand{\sY}{\mathscr{Y}}
\newcommand{\sE}{\mathscr{E}}
\newcommand{\sX}{\mathscr{X}}
\newcommand{\sU}{\mathscr{U}}
\newcommand{\sN}{\mathscr{N}}
\newcommand{\sC}{\mathscr{C}}

\newcommand{\emm}{m}
\newcommand{\ess}{s}

% GHP related symbolds
%\newcommand{\GHPwt}{\circeq}
% \newcommand{\GHPLie}{\text{\L}}
%
%\newcommand{\sib}{\bar{\sigma}}
%\newcommand{\epb}{\bar{\epsilon}}
%\newcommand{\kab}{\bar{\kappa}}
%\newcommand{\beb}{\bar{\beta}}
% \newcommand{\rhb}{\bar{\rho}}
%\newcommand{\pib}{\bar{\pi}}
%\newcommand{\alb}{\bar{\alpha}}
% \newcommand{\tab}{\bar{\tau}}
%\newcommand{\gab}{\bar{\gamma}}
%\newcommand{\mub}{\bar{\mu}}
%\newcommand{\lab}{\bar{\lambda}}
%\newcommand{\nub}{\bar{\nu}}
%
% \newcommand{\kap}{\kappa'}
% \newcommand{\sip}{\sigma'}
% \newcommand{\rhp}{\rho'}
% \newcommand{\tap}{\tau'}
%\newcommand{\bep}{\beta'}
%\newcommand{\epp}{\epsilon'}

\font\ec=ecrm0800 at 12pt
\newcommand{\thorn}{\hbox{\ec\char'336}}
\renewcommand{\eth}{\hbox{\ec\char'360}}
%\renewcommand{\th}{\hbox{\ec\char'336}}
%\newcommand{\edth}{\hbox{\ec\char'360}}
%\newcommand{\thp}{\hbox{\ec\char'336}'}
%\newcommand{\edthp}{\hbox{\ec\char'360}'}
%\newcommand{\Dbar}{\tilde \edth}
%\newcommand{\Nbar}{\tilde \th'}
%\newcommand{\mb}{\bar{m}}

\newcommand{\pheq}{\phantom{=}}

% Chandrasekhar ops
\newcommand{\chand}[2]{\,{}_{#1} \mathcal{L}_{#2}}

% Theorems, etc.
\newtheorem{theorem}{Theorem}
\newtheorem*{corollary}{Corollary}
\newtheorem{lemma}[theorem]{Lemma}
\newtheorem{definition}{Definition}
\newtheorem{proposition}{Proposition}
\newtheorem{exercise}{Exercise}


% Real and imaginary part symbols
\renewcommand{\Re}{\operatorname{Re}}
\renewcommand{\Im}{\operatorname{Im}}

% Spheroidal harmonics
\newcommand{\swsY}[3]{\,{}_{#1}Y_{#2\, #3}}
\newcommand{\sws}[2]{\,{}_{#1}#2}
%\newcommand{\swss}[2]{S_{\emm, \lambda}^{[#1], (#2)}}
%\newcommand{\swsy}[1]{\, {}_{#1}Y^{\phantom{\textrm{p}}}_{\ell, \emm}}
%\newcommand{\swsl}[2]{\lambda^{[#1], (#2)}_\emm}
%\newcommand{\swsc}[2]{C_{\emm, \lambda}^{[#1], (#2)}}
%\newcommand{\swsp}[2]{P_{\emm, \lambda}^{[#1], (#2)}}

%% Sph. harmonic defs (Peter)
%\newcommand{\Y}[1]{\, {}_{#1}Y^{\phantom{\textrm{p}}}_{\ell, \emm}}
%\newcommand{\bY}[1]{\, {}_{#1}\bar{Y}^{\phantom{\textrm{p}}}_{\ell, \emm}}
%\newcommand{\Yp}[1]{\, {}_{#1}Y^\textrm{p}_{\ell, \emm}}
%\newcommand{\bYp}[1]{\, {}_{#1}\bar{Y}^\textrm{p}_{\ell, \emm}}
%\newcommand{\Ym}[2]{\, {}_{#1}Y_{\ell, #2}}
%\newcommand{\Ylm}[3]{\, {}_{#1}Y_{#2, #3}}
%\newcommand{\bYlm}[3]{\, {}_{#1} \bar{Y}_{#2, #3}}
%\newcommand{\C}[1]{\lambda_{\ell, #1}}
%\newcommand{\blambda}{{\mkern0.75mu\mathchar '26\mkern -9.75mu\lambda}}

% Other mathematical symbols
\newcommand{\widebar}[1]{\mathop{\overline{#1}}}
\newcommand{\preindex}[2]{\,{}_{#2}#1}
\newcommand{\half}{\tfrac{1}{2}}
\newcommand{\abs}[1]{\left\lvert #1 \right\rvert}
\newcommand{\norm}[1]{\left\lVert #1 \right\rVert}
% \newcommand{\pd}{\partial}
\newcommand{\dbd}[1]{\dfrac{\partial}{\partial #1}}
\newcommand{\intd}[1]{\textrm{d} #1}
\newcommand{\lie}[1]{\mathcal{L}_{#1}}
% \newcommand{\covlie}[1]{\L{}_#1}
\newcommand{\sforall}{\quad \forall}

% Sections and Subsections
\renewcommand{\thesection}{\Roman{section}} 
\renewcommand{\thesubsection}{\arabic{subsection}}
\renewcommand{\thesubsubsection}{\thesubsection.\arabic{subsubsection}}

%% Equations
%%\def\ben{\begin{equation}}
%%\def\een{\end{equation}}
%%\def\bena{\begin{eqnarray}}
%%\def\eena{\end{eqnarray}}
%\newcommand{\non}{\nonumber}

% Standard operators
\renewcommand{\Re}{\operatorname{Re}}
\renewcommand{\Im}{\operatorname{Im}}

% Differential operators
\newcommand{\dd}{\textrm{d}}
\newcommand{\cX}{\mathcal{X}}

% Spaces
\newcommand{\RR}{\mathbb{R}}
\newcommand{\CC}{\mathbb{C}}

% Vahid's commands
\newcommand{\GHPw}[2]{\left\{ #1, #2 \right\}}
\newcommand{\nls}{\phantom{\frac{1}{\rho}}}

% Christiane's commands
\newcommand{\rC}{\mathrm{C}}
\newcommand{\rN}{\mathrm{N}}
\newcommand{\rU}{\mathrm{U}}
\newcommand{\rI}{\mathrm{I}}
\newcommand{\rII}{\mathrm{II}}
\newcommand{\rIII}{\mathrm{III}}
\newcommand{\rIV}{\mathrm{IV}}

\newcommand{\mA}{\mathcal{A}}
\newcommand{\mB}{\mathcal{B}}
\newcommand{\mC}{\mathcal{C}}
\newcommand{\mD}{\mathcal{D}}
\newcommand{\mE}{\mathcal{E}}
\newcommand{\mF}{\mathcal{F}}
\newcommand{\mH}{\mathcal{H}}
\newcommand{\scrI}{\mathcal{I}}
\newcommand{\mK}{\mathcal{K}}
\newcommand{\mL}{\mathcal{L}}
\newcommand{\mM}{\mathcal{M}}
\newcommand{\mN}{\mathcal{N}}
\newcommand{\mO}{\mathcal{O}}
\newcommand{\mR}{\mathcal{R}}
\newcommand{\mS}{\mathcal{S}}
\newcommand{\mT}{\mathcal{T}}
\newcommand{\mU}{\mathcal{U}}
\newcommand{\mV}{\mathcal{V}}

\newcommand{\bC}{\mathbb{C}}
\newcommand{\bD}{\mathbb{D}}
\newcommand{\bN}{\mathbb{N}}
\newcommand{\bP}{\mathbb{P}}
\newcommand{\bQ}{\mathbb{Q}}
\newcommand{\bR}{\mathbb{R}}
\newcommand{\bS}{\mathbb{S}}
\newcommand{\bZ}{\mathbb{Z}}

\newcommand{\fL}{\mathfrak{L}}
\newcommand{\fM}{\mathfrak{M}}
\newcommand{\fN}{\mathfrak{N}}

\newcommand{\mHc}{\mathcal{H}_c}
\newcommand{\mCH}{\mathcal{CH}}
\newcommand{\ome}{\omega_{\rI}}
\newcommand{\omi}{\omega_{\rII}}
\newcommand{\ri}{\mathrm{in}}
\newcommand{\ru}{\mathrm{up}}
\newcommand{\omegap}{\omega^+}

\DeclareMathOperator{\WF}{WF}
\DeclareMathOperator{\supp}{supp}
\DeclareMathOperator{\csch}{csch}
\DeclareMathOperator{\tr}{tr}
\DeclareMathOperator{\id}{id}
\DeclareMathOperator{\sign}{sign}
\DeclareMathOperator{\Diff}{Diff}

\DeclareOldFontCommand{\bf}{\normalfont\bfseries}{\mathbf}
\DeclareOldFontCommand{\rm}{\normalfont\bfseries}{\mathrm}

% Udefineds
\newcommand{\dOm}{\textrm{UNDEFINED}}
\newcommand{\vt}[1]{\textcolor{red}{Vahid: #1}}
\newcommand{\sh}[1]{\textcolor{red}{Stefan: #1}} 


\newcommand{\ud}{d}
\newcommand{\del}{\partial}


\newcommand{\eps}{\varepsilon}
\newcommand{\vth}{{\vartheta}}

\newcommand{\rin}{{\mathrm{in}}}
\newcommand{\rup}{{\mathrm{up}}}
\newcommand{\rout}{{\mathrm{out}}}
\newcommand{\rdown}{{\mathrm{down}}}
\newcommand{\reg}{{\mathrm{reg}}}
\newcommand{\sing}{{\mathrm{sing}}}

\DeclareMathOperator{\Real}{Re}
\DeclareMathOperator{\Imag}{Im}
\DeclareMathOperator{\Tr}{Tr}
\DeclareMathOperator{\diag}{diag}

\newcommand{\defeq}{\mathrel{:=}}
\newcommand{\F}{{}_2 F_1}
\newcommand{\btheta}{\boldsymbol{\theta}}
\newcommand{\bsigma}{\boldsymbol{\sigma}}
\renewcommand{\bm}{\boldsymbol{m}}
\newcommand{\nn}{\nonumber}

%Peter's shortcuts
\newcommand{\bphi}{\bar \phi}
\newcommand{\bx}{\bar x}
\newcommand{\bt}{\bar t}


\begin{document}

\title{Fist Order Metric Reconstruction in Schwarzschld spacetime}


\author{Vahid Toomani}
\email{vahid.toomani.2@uni-leipzig.de}
\affiliation{Institut f\"ur Theoretische Physik, Universit\"at Leipzig, Leipzig, Germany and \\
Max-Planck-Institute MiS, Leipzig, Germany}

%\author{Stefan Hollands}
%\email{stefan.hollands@uni-leipzig.de}
%\affiliation{Institut f\"ur Theoretische Physik, Universit\"at Leipzig, Leipzig, Germany and \\
%Max-Planck-Institute MiS, Leipzig, Germany}

\date{\today}


\begin{abstract}
My notes!
\end{abstract}

\maketitle
\tableofcontents
%\newpage

%%%%%%%%%%%%%%%%%%%%%%%%%%%%%%%%%%%%%%%%%%%%%%%%%

\section{GHP \& Held}

Combining eq. 123, eq. 124 and eq. J.1a of \cite{Hollands:2024iqp} gives

\begin{equation}
\zeta = - \Psi^{-\sfrac{1}{6}}_2 \bar{\Psi}^{-\sfrac{1}{6}}_2 \rho^{-\sfrac{1}{2}} \bar{\rho}^{\sfrac{1}{2}}
\end{equation}

\begin{equation}
\begin{split}
B_a &= -\rho n_a + \tau \bar{m}_a,\\
B'_a &= -\rho' l_a + \tau' m_a.
\end{split}
\end{equation}

\begin{equation}
t_a = \zeta (B_a - B'_a )
\end{equation}

\begin{equation}
\L_t \eta = t^a \Theta_a \eta - \frac{p}{2} \zeta \Psi_2  \eta - \frac{q}{2} \bar{\zeta} \bar{\Psi}_2  \eta = \zeta \rho' \thorn \eta -\zeta \rho \thorn' \eta - \zeta \tau' \eth \eta + \zeta \tau \eth'' \eta - \frac{p}{2} \zeta \Psi_2  \eta - \frac{q}{2} \bar{\zeta} \bar{\Psi}_2  \eta
\end{equation}


\begin{equation}
\begin{split}
\thorn \thorn \thorn \thorn \left( \zeta^4 \psi_4 \right)&= \eth' \eth' \eth' \eth' \left( \zeta^4 \psi_0 \right)  - 3 \L_t \bar{\psi}_0,\\
\thorn' \thorn' \thorn' \thorn' \left( \zeta^4 \psi_0 \right)&= \eth \eth \eth \eth \left( \zeta^4 \psi_4 \right)  + 3 \L_t \bar{\psi}_4.
\end{split}
\end{equation}

\section{Components of Homogenous $\Phi$}

In this section, I express how one achieves a metric perturbation that falls fast enough at $\mathscr{I}^+$. To achieve this end I used the procedure demonstrated in details by Hollands et. al. in \cite{Hollands:2024iqp}.
This method is based on the GHZ reconstruction \cite{Green:2019nam} combined with the puncture scheme \cite{Toomani:2021jlo}.
The puncture scheme implemented in \cite{Toomani:2021jlo} breaks the retarded solution $h^\text{ret}_{ab}$ into two parts of punctured field $h^\text{P}_{ab}$ which contains the singular field field and the residual part $h^\text{Res}_{ab}$. The residual field satisfy a homogeneous linearized Einstein equation sourced by $\mathcal{E} h^\text{P}_{ab}$ away from the singular source, i.e.,

\begin{equation}\label{LEEofResField}
\mathcal{E} h^\text{Res}_{ab} = - \mathcal{E} h^\text{P}_{ab}.
\end{equation}

The puncture field is an estimate for the singular field in a neighborhood of the orbiting compact object. This estimate is given by the product of the singular field and a $W(r, \theta, \phi)$ with the following properties:

\begin{enumerate}
\item $W$ is compactly supported (usually $W = 0$ for $r < r_\text{min}$ and $r > r_\text{max}$, the support of $W$ is called puncture region;
\item $W = 1$ on the location of the compact object.
\end{enumerate}

One obvious choice would be $W(r) = \Theta(r - r_\text{min}) - \Theta(r_\text{max} - r)$ but one might choose a smooth one too. Consequently, residual field $h^\text{Res}_{ab}$ is only sourced in the puncture region.

Applying $\mathcal{S}$ on both sides of \eqref{LEEofResField} and using the Wald identity \cite{Wald identity paper},

\begin{equation}
\mathcal{O} \psi^\text{Res}_0 = - \mathcal{S}^{ab} T^\text{P}_{ab} = - T^\text{P}_0.
\end{equation}

Using the Wald identity \cite{Wald identity paper}, we have



Combining eq. 123, eq. 124 and eq. J.1a of \cite{Hollands:2024iqp} gives

\begin{equation}
2 \tilde{\thorn}' \tilde{\thorn}' \tilde{\thorn}' \bar{\psi}^{5 \circ}_0 = \tilde{\eth}' \tilde{\eth}' \tilde{\eth}' \tilde{\eth}' \tilde{\eth} \tilde{\eth} \tilde{\eth} \tilde{\eth} \bar{f}^\circ - 9 \Psi^\circ \bar{\Psi}^\circ \tilde{\thorn}' \tilde{\thorn}' \bar{f}^\circ.
\end{equation}

By expanding $\bar{f}^\circ$ and $\bar{\psi}^{5 \circ}_0$ in spheroidal (spin-weighted harmonics in Schwarzschild) harmonics, one can reverse this equation and obtain $\bar{f}^\circ$. Next, we obtain foloowing Held-like scalars $\Phi^{i \circ}$ from $\bar{f}^\circ$ (and potentially ODEs in terms of $\frac{\partial}{\partial u}$).

\begin{equation}
\tilde{\thorn}' \tilde{\thorn}' \Phi^{0 \circ} = - \frac{1}{3} \tilde{\eth}' \tilde{\eth}' \tilde{\eth}' \tilde{\eth} \tilde{\eth} \tilde{\eth} \bar{f}^\circ - \left[ 2 \bar{\tau}^\circ \tilde{\eth}' \tilde{\eth} \tilde{\eth} + \Psi^\circ \tilde{\eth}' \tilde{\eth} + 2 \Psi^\circ \rho^{\prime \circ} \right] \tilde{\thorn}' \bar{f}^\circ
\end{equation}

\begin{equation}
\tilde{\thorn}' \Phi^{1 \circ} = - \tilde{\eth}' \tilde{\eth}' \tilde{\eth} \tilde{\eth} \bar{f}^\circ - \left[ 4 \bar{\tau}^\circ \tilde{\eth} + 3 \Psi^\circ \right] \tilde{\thorn}' \bar{f}^\circ
\end{equation}

\begin{equation}
\Phi^{2 \circ} = -2 \tilde{\eth}' \tilde{\eth} \bar{f}^\circ,
\end{equation}

\begin{equation}
\Phi^{3 \circ} = -2 \tilde{\thorn}' \bar{f}^\circ,
\end{equation}

\begin{equation}
\bar{\zeta}^\circ_m = \tilde{\eth} \bar{f}^\circ,
\end{equation}

\begin{equation}
\tilde{\thorn}' \zeta^\circ_l = - \Re \left\{ \tilde{\eth} \tilde{\eth} \bar{f}^\circ \right\},
\end{equation}

\begin{equation}
\tilde{\thorn}' \zeta^\circ_n = - \Re \left\{ \tilde{\eth}' \tilde{\eth} \tilde{\eth} \tilde{\eth} \bar{f}^\circ - \rho^{\prime \circ} \tilde{\eth} \tilde{\eth} \bar{f}^\circ - 2 \tau^\circ \tilde{\eth} \tilde{\thorn}' \bar{f}^\circ - \frac{1}{2} \Omega^\circ \tilde{\eth} \tilde{\eth} \tilde{\thorn}' \bar{f}^\circ \right\},
\end{equation}

kslfhadlf

\begin{equation}
\begin{split}
h_{ab} & \to h_{ab} - \mathcal{L}_\zeta g_{ab}\\
\Phi & \to \Phi + \Phi^{0 \circ} + \Phi^{1 \circ} \frac{1}{\rho} + \Phi^{2 \circ} \frac{1}{\rho^2} + \Phi^{3 \circ} \frac{1}{\rho^3}
\end{split}
\end{equation}

\begin{equation}
\begin{split}
\mathcal{L}_\zeta g_{ab} l^a & = 0\\
\mathcal{L}_\zeta g_{ab} m^a \bar{m}^b & = 0
\end{split}
\end{equation}

\begin{multline}
\mathcal{L}_\zeta g_{ab} n^a n^b = \Re \bigg[ 4 \rho \bar{\rho} \bar{\tau}^\circ \tilde{\eth}' \tilde{\eth} \tilde{\eth} \zeta^\circ_l - 4 \rho^2 \bar{\rho} \bar{\tau}^{\circ 2} \tilde{\eth} \tilde{\eth} \zeta^\circ_l + 2 \rho^2 \Psi^\circ \tilde{\eth}' \tilde{\eth} \zeta^\circ_l + 4 \rho^2 \bar{\rho} \Psi^\circ \bar{\tau}^\circ \tilde{\eth} \zeta^\circ_l\\
- 2 \tilde{\eth}' \tilde{\eth}' \tilde{\eth}' \tilde{\eth} f^\circ - 2\left( \frac{1}{\rho} + 4 \Omega^\circ \right) \tilde{\eth}' \tilde{\eth}' \tilde{\thorn}' f^\circ - 4 \rho \bar{\rho} \Omega^\circ \bar{\tau}^\circ \tilde{\eth}' \tilde{\eth}' \tilde{\eth} f^\circ + 4 \left( 4 + \bar{\rho} \Omega^\circ - \rho \bar{\rho} \Omega^{\circ 2} \right) \bar{\tau}^\circ \tilde{\eth}' \tilde{\thorn}' f^\circ\\
- \left( 8 \rho^{\prime \circ} + 3 \rho \Psi^\circ + 3 \bar{\rho} \bar{\Psi}^\circ + 2 \rho^2 \Omega^{\circ 2} \Psi^\circ + 8 \rho \bar{\rho} \tau^\circ \bar{\tau}^\circ \right) \tilde{\eth}' \tilde{\eth}' f^\circ + 4 \rho^2 \bar{\tau}^{\circ 2} \tilde{\eth}' \tilde{\eth} f^\circ\\
+ 4 \rho \left( \rho + 2 \bar{\rho} \right) \Omega^\circ \bar{\tau}^{\circ 2} \tilde{\thorn}' f^\circ + 8 \rho^2 \bar{\rho} \tau^\circ \bar{\tau}^{\circ 2} \tilde{\eth}' f^\circ + 4 \rho \bar{\rho} \bar{\Psi}^\circ \bar{\tau}^\circ \tilde{\eth}' f^\circ - 24 \rho \bar{\rho} \Omega^\circ \rho^{\prime \circ} \bar{\tau}^\circ \tilde{\eth}' f^\circ + 16 \rho^2 \rho^{\prime \circ} \bar{\tau}^{\circ 2} f^\circ \bigg]
\end{multline}

\begin{multline}
\mathcal{L}_\zeta g_{ab} n^a m^b = \bar{\rho} \tilde{\eth}' \tilde{\eth} \tilde{\eth} \zeta^\circ_l - 2 \rho \bar{\rho} \bar{\tau} \tilde{\eth} \tilde{\eth} \zeta^\circ_l + \rho (\rho + \bar{\rho}) \Psi^\circ \tilde{\eth} \zeta^\circ_l - \bar{\rho} \Omega^\circ \tilde{\eth}' \tilde{\eth}' \tilde{\eth} f^\circ + 2 \rho \bar{\tau}^\circ \tilde{\eth}' \tilde{\eth} f^\circ\\
+ \left(\frac{1}{\bar{\rho}} - 2 \bar{\rho} \Omega^{\circ 2} \right) \tilde{\eth}' \tilde{\thorn}' f^\circ + 2 (\rho + \bar{\rho}) \Omega^\circ \bar{\tau}^\circ \tilde{\thorn}' f^\circ - 2 \bar{\rho} \tau^\circ \tilde{\eth}' \tilde{\eth}' f^\circ - (\rho^2 \Omega^\circ + \rho) \Psi^\circ \tilde{\eth}' f^\circ\\
+ \bar{\rho} \bar{\Psi}^\circ \tilde{\eth}' f^\circ - 6 \bar{\rho} \Omega^\circ \rho^{\prime \circ} \tilde{\eth}' f^\circ + 4 \rho \bar{\rho} \tau^\circ \bar{\tau}^\circ \tilde{\eth}' f^\circ + 8 \rho \rho^{\prime \circ} \bar{\tau}^\circ f^\circ
\end{multline}

\begin{equation}
\mathcal{L}_\zeta g_{ab} m^a m^b = - 2 \bar{\rho} \tilde{\eth} \tilde{\eth} \zeta^\circ_l + 2 \tilde{\eth}' \tilde{\eth} f^\circ + 2 \Omega^\circ \tilde{\thorn}' f^\circ + 4 \bar{\rho} \tau^\circ \tilde{\eth}' f^\circ + 8 \rho^{\prime \circ} f^\circ
\end{equation}

\section{Sphroidal Expansions}

\begin{equation}
\psi^{1 \circ}_4 = \sum_{\ell \emm} \int \intd{\omega} \psi^{1 \circ}_{4, \ell \emm}(\omega) \sws{-2}{S_{\ell \emm}} (\theta, \varphi_\star ; a \omega) e^{-i \omega u}
\end{equation}

\begin{equation}
\psi^{5 \circ}_0 = \sum_{\ell \emm} \int \intd{\omega} \psi^{5 \circ}_{0, \ell \emm}(\omega) \sws{+2}{S_{\ell \emm}} (\theta, \varphi_\star ; a \omega) e^{-i \omega u}
\end{equation}

\begin{multline}
\psi^{1 \circ}_{4, \ell \emm}(\omega) = \frac{\omega^4}{\sws{-2}{B_{\ell \emm}}(a \omega) \sws{-2}{\bar{B}_{\ell -\emm}}(-a \omega) + 9 M^2 \omega^2}\\
\left[ \sws{-2}{\bar{B}_{\ell -\emm}}(-a \omega) \psi^{5 \circ}_{0, \ell \emm}(\omega) + 3 i \omega (-1)^\emm \bar{\psi}^{5 \circ}_{0, \ell -\emm}(-\omega) \right]
\end{multline}

\begin{multline}
\bar{h}^{1 \circ}_{m m} = \sum_{\ell \emm} \int \intd{\omega} \frac{-2 \omega^2}{\sws{-2}{B_{\ell \emm}}(a \omega) \sws{-2}{\bar{B}_{\ell -\emm}}(-a \omega) + 9 M^2 \omega^2} \times\\
\left[ \sws{-2}{\bar{B}_{\ell -\emm}}(-a \omega) \psi^{5 \circ}_{0, \ell \emm}(\omega) + 3 i \omega (-1)^\emm \bar{\psi}^{5 \circ}_{0, \ell -\emm}(-\omega) \right] \sws{+2}{S_{\ell \emm}} (\theta, \varphi_\star ; a \omega) e^{-i \omega u}
\end{multline}

\begin{equation}
f^\circ = \sum_{\ell \emm} \int \intd{\omega} \frac{2 i \omega^3}{\sws{-2}{B_{\ell \emm}}(a \omega) \sws{-2}{\bar{B}_{\ell -\emm}}(-a \omega) + 9 M^2 \omega^2} \psi^{5 \circ}_{0, \ell \emm}(\omega) \sws{+2}{S_{\ell \emm}} (\theta, \varphi_\star ; a \omega) e^{-i \omega u}
\end{equation}

\begin{equation}
e^\circ = \sum_{\ell \emm} \int \intd{\omega} \frac{-4 \omega^4 (-1)^\emm}{\sws{-2}{B_{\ell \emm}}(a \omega) \sws{-2}{\bar{B}_{\ell -\emm}}(-a \omega) + 9 M^2 \omega^2} \bar{\psi}^{5 \circ}_{0, \ell -\emm}(-\omega) \sws{-2}{S_{\ell \emm}} (\theta, \varphi_\star ; a \omega) e^{-i \omega u}
\end{equation}

\section{Schwarzschild Case}

\begin{equation}
\sws{-2}{B_{\ell -\emm}}(a \omega) = \sws{2}{B_\ell} = \frac{1}{4} \frac{(\ell + 2)!}{(\ell - 2)!} = \frac{1}{4} (\ell +2) (\ell + 1) \ell (\ell - 1)
\end{equation}

\begin{equation}
\sws{-1}{B_{\ell -\emm}}(a \omega) = \sws{1}{B_\ell} = \frac{1}{2} \frac{(\ell + 1)!}{(\ell - 1)!} = \frac{1}{2} \ell (\ell + 1)
\end{equation}

\begin{equation}
\Phi^{0 \circ} = d^\circ = \sum_{\ell \emm} \int \intd{\omega} \frac{ - \sfrac{2}{3} i \omega (-1)^\emm}{\sws{2}{B^2_\ell} + 9 M^2 \omega^2} \left[ \sws{2}{B_\ell} \sws{1}{B_\ell} - 3iM \omega (1 + \sqrt{\sws{2}{B_{\ell}}}) \right] \bar{\psi}^{5 \circ}_{0, \ell -\emm}(- \omega) \sws{-2}{Y_{\ell \emm}} (\theta, \varphi_\star) e^{-i \omega u}
\end{equation}

\begin{equation}
\Phi^{1 \circ} = \sum_{\ell \emm} \int \intd{\omega} \frac{ 2 \omega^2 (-1)^\emm}{\sws{2}{B^2_\ell} + 9 M^2 \omega^2} \left[ \sws{2}{B_\ell} - 3iM \omega \right] \bar{\psi}^{5 \circ}_{0, \ell -\emm}(- \omega) \sws{-2}{Y_{\ell \emm}} (\theta, \varphi_\star) e^{-i \omega u}
\end{equation}

\begin{equation}
\Phi^{2 \circ} = \sum_{\ell \emm} \int \intd{\omega} \frac{ 4 i \omega^3 \sfrac{\sws{2}{B_\ell}}{\sws{1}{B_\ell}} (-1)^\emm}{\sws{2}{B^2_\ell} + 9 M^2 \omega^2} \bar{\psi}^{5 \circ}_{0, \ell -\emm}(- \omega) \sws{-2}{Y_{\ell \emm}} (\theta, \varphi_\star) e^{-i \omega u}
\end{equation}

\begin{equation}
\Phi^{3 \circ} = e^\circ = \sum_{\ell \emm} \int \intd{\omega} \frac{-4 \omega^4 (-1)^\emm}{\sws{2}{B^2_\ell} + 9 M^2 \omega^2} \bar{\psi}^{5 \circ}_{0, \ell -\emm}(-\omega) \sws{-2}{Y_{\ell \emm}} (\theta, \varphi_\star) e^{-i \omega u}
\end{equation}

\begin{equation}
\zeta^\circ_m = \sum_{\ell \emm} \int \intd{\omega} \frac{-2 i \omega^3 \sqrt{\sfrac{\sws{2}{B_\ell}}{\sws{1}{B_\ell}}}}{\sws{2}{B^2_\ell} + 9 M^2 \omega^2} \psi^{5 \circ}_{0, \ell \emm}(\omega) \sws{+1}{Y_{\ell \emm}} (\theta, \varphi_\star) e^{-i \omega u}
\end{equation}

\begin{equation}
\zeta^\circ_l = \sum_{\ell \emm} \int \intd{\omega} \frac{\omega^2 \sqrt{\sws{2}{B_\ell}}}{\sws{2}{B^2_\ell} + 9 M^2 \omega^2} \psi^{5 \circ}_{0, \ell \emm}(\omega) \sws{0}{Y_{\ell \emm}} (\theta, \varphi_\star) e^{-i \omega u} + cc.
\end{equation}

\begin{equation}
\zeta^\circ_n = \sum_{\ell \emm} \int \intd{\omega} \frac{-\omega^2 \sqrt{\sws{2}{B_\ell}}}{\sws{2}{B^2_\ell} + 9 M^2 \omega^2} \left( \sws{1}{B_\ell} + \frac{1}{2} \right)\psi^{5 \circ}_{0, \ell \emm}(\omega) \sws{0}{Y_{\ell \emm}} (\theta, \varphi_\star) e^{-i \omega u} + cc.
\end{equation}

%%%%%%%%%%%%%%%%%%%%%%%%%%%%%%%%%%%%%%%%%%%%%%%%%


\appendix

%%%%%%%%%%%%%%%%%%%%%%%%%%%%%%%%%%%%%%%%%%%%%%%%%

%\include{section 1.tex}

%%%%%%%%%%%%%%%%%%%%%%%%%%%%%%%%%%%%%%%%%%%%%%%%%

\bibliography{references.bib}


\end{document}

%@article{PhysRevA.42.1027,
%  title = {Gaussian-orthogonal-ensemble level statistics in a one-dimensional system},
%  author = {Wu, Hua and Valli\`eres, Michel and Feng, Da Hsuan and Sprung, Donald W. L.},
%  journal = {Phys. Rev. A},
%  volume = {42},
%  issue = {3},
%  pages = {1027--1032},
%  numpages = {0},
%  year = {1990},
%  month = {Aug},
%  publisher = {American Physical Society},
%  doi = {10.1103/PhysRevA.42.1027},
%  url = {https://link.aps.org/doi/10.1103/PhysRevA.42.1027}
%}
%
%%
